\begin{tikzpicture}[x=1cm,y=1cm,>=stealth,thick]
  \tikzstyle{layer}=[draw, rounded corners=5pt, minimum width=13.5cm, minimum height=1.0cm, align=center]
  \tikzstyle{side}=[draw, rounded corners=4pt, minimum width=3.5cm, minimum height=0.8cm, align=center, font=\small]

  \node[layer, fill=blue!8] (ui) at (0,0) {\textbf{Танилцуулгын давхарга}\\Веб хуудас, компонент, хэрэглэгчийн интерфейс, курс/хичээл/чат};
  \node[layer, fill=green!8] (app) at (0,-1.8) {\textbf{Аппликейшний давхарга}\\Клиент төлөвийн функц, хандалтын хяналт, хичээлийн урсгал, сорилын төлөв};
  \node[layer, fill=yellow!10] (api) at (0,-3.6) {\textbf{Хиймэл оюуны туслах давхарга}\\Програмчлалын интерфейсийн зам, сэргээн хайлт, зааварчилгаа, нөөц хариу};
  \node[layer, fill=orange!10] (data) at (0,-5.4) {\textbf{Өгөгдлийн давхарга}\\Үүлэн харилцаат өгөгдлийн сан, нэвтрэлт, мөрийн түвшний хамгаалалт, ахиц};
  \node[layer, fill=gray!14] (assets) at (0,-7.2) {\textbf{Контент ба өгөгдлийн багцын давхарга}\\Курсийн эх өгөгдөл, видео хичээлийн лавлагаа, өгөгдлийн багц, ФЛ Студио зөвлөмж};

  \draw[->] (ui) -- (app);
  \draw[->] (app) -- (api);
  \draw[->] (api) -- (data);
  \draw[->] (data) -- (assets);

  \node[side, fill=purple!8] (student) at (-9.0,-0.9) {Суралцагч};
  \node[side, fill=purple!8] (admin) at (9.0,-0.9) {Админ};
  \node[side, fill=cyan!8] (llm) at (9.0,-3.6) {Хэлний загварын үйлчилгээ};
  \node[side, fill=cyan!8] (supabase) at (9.0,-5.4) {Үүлэн сангийн үйлчилгээ};

  \draw[->] (student) -- (ui.west);
  \draw[->] (admin) -- (ui.east);
  \draw[->] (api.east) -- (llm);
  \draw[->] (data.east) -- (supabase);
\end{tikzpicture}
